% =====================================================================
\chapter{Espacios de probabilidad}
\label{tema:3}
% =====================================================================

Este tema da la definición axiomática de probabilidad, deduce de los axiomas las reglas de cálculo y estudia el caso más importante para empezar, el modelo de Laplace, en el que calcular una probabilidad se reduce a contar. Se siguen la sección §1.2 de Bertsekas y Tsitsiklis~\cite{bt} (clase~1 de MIT 6.041~\cite{mit6041f10}), el capítulo~2 de Ross~\cite{ross} (§2.2 a §2.5), las secciones §1.2, §1.3 y §1.6 de Blitzstein y Hwang~\cite{bh} y el capítulo~16, «Events and Probability Spaces», de~\cite{mcs}, que incluye el problema del cumpleaños (§16.4).

\section{Experimentos, espacio muestral y sucesos}

\begin{definicion}[Espacio muestral y sucesos]
Un \emph{experimento aleatorio} es un procedimiento cuyo resultado no puede predecirse con certeza. El \emph{espacio muestral} $\Omega$ es el conjunto de todos sus resultados posibles. Un \emph{suceso} es un subconjunto $A\subseteq\Omega$; se dice que $A$ \emph{ocurre} si el resultado obtenido pertenece a $A$. Los sucesos $\{\omega\}$ se llaman \emph{elementales}.
\end{definicion}

Las operaciones con conjuntos del módulo~0.2 se traducen en operaciones con sucesos: $A\cup B$ es «ocurre $A$ o $B$»; $A\cap B$, «ocurren $A$ y $B$»; $A^{c}=\Omega\setminus A$, «no ocurre $A$»; $A\subseteq B$, «si ocurre $A$, ocurre $B$». Dos sucesos son \emph{incompatibles} (o mutuamente excluyentes) si $A\cap B=\emptyset$. Las leyes de De Morgan dan, por ejemplo, $(A\cup B)^{c}=A^{c}\cap B^{c}$: «no ocurre ninguno de los dos».

\addtocounter{definicion}{1}%

\section{Axiomas de Kolmogorov}

\begin{definicion}[Espacio de probabilidad]\label{def:t3-axiomas}
Una \emph{ley de probabilidad} sobre $\Omega$ es una función $\PP$ que asigna a cada suceso $A$ un número real $\PP(A)$ y que cumple:
\begin{enumerate}
  \item \emph{no negatividad}: $\PP(A)\ge0$ para todo suceso $A$;
  \item \emph{normalización}: $\PP(\Omega)=1$;
  \item \emph{aditividad numerable}: si $A_1,A_2,\dots$ son sucesos incompatibles dos a dos, $\PP\bigl(\bigcup_{i\ge1}A_i\bigr)=\sum_{i\ge1}\PP(A_i)$.
\end{enumerate}
El par $(\Omega,\PP)$ se llama \emph{espacio de probabilidad}.
\end{definicion}

\begin{observacion}
Si $\Omega$ es finito o numerable, todo subconjunto puede ser un suceso. Si $\Omega$ es no numerable, como un intervalo de $\R$, no es posible en general asignar probabilidad a todos los subconjuntos de forma coherente, y los sucesos se restringen a una familia adecuada (una $\sigma$-álgebra). En este módulo esa sutileza no interviene: todos los sucesos que aparecen (intervalos, uniones de intervalos, regiones del plano) son admisibles~\cite[§1.6]{bh}.
\end{observacion}

\section{Propiedades que se deducen de los axiomas}

\begin{proposicion}\label{prop:t3-props}
Sean $A$, $B$ y $A_1,\dots,A_n$ sucesos. Entonces:
\begin{enumerate}
  \item $\PP(\emptyset)=0$;
  \item \emph{aditividad finita}: si $A_1,\dots,A_n$ son incompatibles dos a dos, $\PP\bigl(\bigcup_{i=1}^{n}A_i\bigr)=\sum_{i=1}^{n}\PP(A_i)$;
  \item $\PP(A^{c})=1-\PP(A)$;
  \item si $A\subseteq B$, entonces $\PP(B\setminus A)=\PP(B)-\PP(A)$ y $\PP(A)\le\PP(B)$;
  \item $0\le\PP(A)\le1$;
  \item $\PP(A\cup B)=\PP(A)+\PP(B)-\PP(A\cap B)$.
\end{enumerate}
\end{proposicion}

\begin{teorema}[Desigualdad de Boole o cota de la unión]\label{teo:t3-boole}
Para sucesos cualesquiera $A_1,\dots,A_n$,
\(
\PP\bigl(\bigcup_{i=1}^{n}A_i\bigr)\le\sum_{i=1}^{n}\PP(A_i).
\)
\end{teorema}

La fórmula de inclusión--exclusión del \cref{teo:t2-ie} tiene una versión probabilística con la misma estructura; para tres sucesos,
\[
\PP(A\cup B\cup C)=\PP(A)+\PP(B)+\PP(C)-\PP(A\cap B)-\PP(A\cap C)-\PP(B\cap C)+\PP(A\cap B\cap C),
\]
y se demuestra a partir de (vi) exactamente como el \cref{teo:t1-ie}.

\section{Espacios discretos y modelo de Laplace}

\begin{proposicion}[Espacios discretos]\label{prop:t3-discreto}
Sea $\Omega$ finito o numerable y $p\colon\Omega\to[0,1]$ con $\sum_{\omega\in\Omega}p(\omega)=1$. Entonces $\PP(A)=\sum_{\omega\in A}p(\omega)$ define una ley de probabilidad. Recíprocamente, toda ley de probabilidad sobre $\Omega$ es de esta forma, con $p(\omega)=\PP(\{\omega\})$.
\end{proposicion}

\begin{definicion}[Modelo de Laplace]\label{def:t3-laplace}
Si $\Omega$ es finito y todos los resultados son equiprobables, $p(\omega)=1/\abs{\Omega}$, y entonces
\[
\PP(A)=\frac{\abs{A}}{\abs{\Omega}}=\frac{\text{casos favorables}}{\text{casos posibles}}.
\]
\end{definicion}

El modelo de Laplace es una \emph{hipótesis} sobre el experimento, no una consecuencia de los axiomas: es razonable para dados, monedas y cartas bien mezcladas, pero no, por ejemplo, para la suma de dos dados, cuyos 11 valores no son equiprobables. La técnica habitual consiste en elegir un espacio muestral en el que la equiprobabilidad sea evidente, normalmente uno de resultados \emph{ordenados} y distinguibles, aunque los objetos del problema parezcan indistinguibles.

\addtocounter{definicion}{1}%

\section{El problema del cumpleaños}

\begin{proposicion}\label{prop:t3-cumple}
Se asigna a cada una de $n$ personas un día entre $N$, de forma uniforme e independiente (todas las $N^{n}$ asignaciones son equiprobables). La probabilidad de que dos o más personas compartan día es
\[
p_n=1-\frac{N^{\underline{n}}}{N^{n}}=1-\prod_{i=0}^{n-1}\Bigl(1-\frac{i}{N}\Bigr),
\]
y además $1-p_n\le\exp\!\bigl(-\tfrac{n(n-1)}{2N}\bigr)$.
\end{proposicion}

\addtocounter{definicion}{1}%

\addtocounter{figure}{1}%

\addtocounter{definicion}{1}%

\addtocounter{definicion}{1}%

